\documentclass[12pt,a4paper]{article}
\usepackage[utf8]{inputenc}
\usepackage{xeCJK}
\usepackage{geometry}
\usepackage{booktabs}
\usepackage{array}
\usepackage{longtable}
\usepackage{multirow}
\usepackage{graphicx}
\usepackage{fancyhdr}
\usepackage{color}
\usepackage{xcolor}
\usepackage{amsmath}
\usepackage{amsfonts}
\usepackage{amssymb}

\geometry{left=2cm,right=2cm,top=2.5cm,bottom=2.5cm}

\setCJKmainfont{Noto Serif CJK TC}
\setCJKsansfont{Noto Sans CJK TC}
\setCJKmonofont{Noto Sans CJK TC}


\pagestyle{fancy}
\fancyhf{}
\rhead{\thepage}
\lhead{心臟內科見習教學經驗總結}

\title{\Large \textbf{心臟內科見習教學經驗總結}\\
\large \textbf{暨計畫書修訂建議}\\
\normalsize Clinical Shadowing Program Experience Summary \& Revision}
\author{謝慕揚 醫師\\
向臨床教師們介紹教學流程與成果}
\date{2025年9月14日}

\begin{document}

\maketitle

\section{教學流程介紹}

\subsection{整體教學架構}

這次兩週的心臟內科見習採用了系統化的教學設計，包含三個核心元素：

\begin{enumerate}
\item \textbf{事前規劃}：詳細的教學計畫書
\item \textbf{過程記錄}：每日學習記錄與分享
\item \textbf{成果評估}：標準化EPA評分系統
\end{enumerate}

\begin{figure}[h]
\centering
\begin{tabular}{|c|c|c|}
\hline
\textbf{Phase 1} & \textbf{Phase 2} & \textbf{Phase 3} \\
事前準備 & 執行階段 & 評估階段 \\
\hline
教學計畫書制定 & Clinical Shadowing & EPA評分 \\
學生預習指導 & 每日記錄分享 & 量化指標統計 \\
目標設定 & 即時指導回饋 & 成果報告撰寫 \\
\hline
\end{tabular}
\caption{三階段教學流程架構}
\end{figure}

\subsection{Phase 1: 事前規劃階段}

\subsubsection{教學計畫書設計}
我們事前制定了完整的《Clinical Shadowing Program in Cardiology》計畫書，包含：

\begin{itemize}
\item 兩週詳細時程表（門診、病房、導管室、超音波室）
\item 預習準備事項（Portal系統、心導管基礎、臨床知識）
\item 技能學習檢核表
\item 安全規範與注意事項
\item 評估方式說明
\end{itemize}

計畫書在見習前一週提供學生詳閱，確保學習目標明確。

\subsubsection{預習要求}
要求學生預習：
\begin{itemize}
\item Portal電腦系統操作
\item 心導管影像判讀基礎
\item 常見心血管疾病診斷標準
\item 心電圖判讀基礎
\end{itemize}

\subsection{Phase 2: Clinical Shadowing執行}

\subsubsection{每日學習安排}
學生每日跟隨我進行：
\begin{itemize}
\item 晨間病房迴診
\item 門診看診（平均每日35-40位病患）
\item 心導管室手術觀摩
\item 超音波室技能訓練
\item 病歷撰寫指導
\end{itemize}

\subsubsection{即時記錄與分享}
學生每日在Google Doc完成學習記錄，格式包含：
\begin{itemize}
\item 今日主要活動
\item 新學到的技能/知識
\item 印象深刻的病例（病例描述、學習重點）
\item 反思思考
\item 今日困難/疑問
\item 明日學習目標
\item 指導醫師回饋
\end{itemize}

透過Google Doc即時分享，我能每日檢視並給予回饋。

\subsection{Phase 3: EPA評估系統}

\subsubsection{評估架構}
我設計了心臟內科專用的10項EPAs評估系統：

\begin{table}[h]
\centering
\caption{心臟內科10項EPAs}
\begin{tabular}{|c|p{8cm}|}
\hline
\textbf{EPA編號} & \textbf{專業活動內容} \\
\hline
EPA 1 & 心臟科病史詢問與風險評估 \\
EPA 2 & 心臟血管系統理學檢查 \\
EPA 3 & 心電圖判讀與分析 \\
EPA 4 & 心臟超音波基礎操作與判讀 \\
EPA 5 & 急性冠心症診斷與初步處理 \\
EPA 6 & 心律不整診斷與處理 \\
EPA 7 & 心衰竭診斷與治療 \\
EPA 8 & 高血壓診斷與管理 \\
EPA 9 & 心導管檢查結果判讀 \\
EPA 10 & 心臟科病患照護協調與溝通 \\
\hline
\end{tabular}
\end{table}

\subsubsection{Level分級標準}
每項EPA評估Level 1-5：
\begin{itemize}
\item Level 1: 僅能觀察學習
\item Level 2: 可在直接監督下執行
\item Level 3: 可獨立執行但需回報指導
\item Level 4: 可獨立執行，監督者隨時可得
\item Level 5: 完全獨立執行並指導他人
\end{itemize}

\section{實際執行成果}

\subsection{量化學習指標}

這次見習的具體量化成果：

\begin{table}[h]
\centering
\caption{周煜珉醫師兩週見習量化成果}
\begin{tabular}{|p{4cm}|p{2cm}|p{6cm}|}
\hline
\textbf{學習項目} & \textbf{完成數量} & \textbf{詳細說明} \\
\hline
病患接診 & 6位 & 完整從入院到出院全程照護 \\
\hline
病歷撰寫 & 6份 & Admission \& Discharge Notes \\
\hline
床邊超音波(POCUS) & 6次 & 心臟基本切面獲取 \\
\hline
動脈血管超音波 & 5次 & 頸動脈、下肢動脈檢查 \\
\hline
靜脈血管超音波 & 5次 & DVT篩檢、下肢靜脈評估 \\
\hline
心導管觀摩 & 10+台 & PCI、PTA、TPM等多種術式 \\
\hline
門診病例觀摩 & 70+例 & 涵蓋主要心血管疾病 \\
\hline
TEE觀摩 & 2次 & 包括LAAO手術前評估 \\
\hline
\end{tabular}
\end{table}

\subsection{EPA評估結果}

最終EPA評估結果：

\begin{table}[h]
\centering
\caption{EPA評估成績分布}
\begin{tabular}{|c|c|c|}
\hline
\textbf{等級} & \textbf{項目數} & \textbf{具體EPAs} \\
\hline
Level 3 & 1項 & EPA 2 (理學檢查) \\
\hline
Level 2 & 4項 & EPA 1, 3, 8, 9 \\
\hline
Level 1 & 3項 & EPA 4, 7, 10 \\
\hline
None & 2項 & EPA 5, 6 (缺乏急症案例) \\
\hline
\end{tabular}
\end{table}

\subsection{最終評分}

採用多元評分系統：
\begin{itemize}
\item EPA基礎能力表現：75\%
\item 學習態度與表現：20\%
\item 特殊貢獻加分：5\%
\end{itemize}

\textcolor{red}{\Large \textbf{最終總分：92分（A級優良）}}

\section{教學心得與觀察}

\subsection{執行要素分析}

\begin{enumerate}
\item \textbf{結構化計畫}：詳細的事前規劃確保學習系統性
\item \textbf{即時記錄}：Google Doc每日分享促進反思與回饋
\item \textbf{量化管理}：具體數字目標提升學習動機
\item \textbf{標準評估}：EPA系統提供客觀評分依據
\end{enumerate}

\subsection{學生表現亮點}

\begin{itemize}
\item \textbf{主動學習}：每日完整記錄，自主閱讀NEJM
\item \textbf{技能積極}：16次超音波操作，6次病患接診
\item \textbf{反思深入}：能整合理論與實務，提出改善想法
\item \textbf{專業態度}：與團隊建立良好關係，獲得肯定
\end{itemize}

\section{計畫書修訂建議}

基於這次實際執行經驗，提出以下修訂建議：

\subsection{1. 建立量化指標系統}

\subsubsection{必達成指標}
在計畫書中明確訂定最低完成標準：

\begin{table}[h]
\centering
\caption{見習期間最低完成指標（修訂建議）}
\begin{tabular}{|p{4cm}|p{2cm}|p{6cm}|}
\hline
\textbf{學習項目} & \textbf{最低要求} & \textbf{評核方式} \\
\hline
病患接診 & 5-8位 & 完整病歷撰寫，指導醫師確認 \\
\hline
Admission Notes & 5-8份 & 格式正確，內容完整 \\
\hline
Discharge Notes & 5-8份 & 包含完整出院計畫與追蹤 \\
\hline
床邊超音波(POCUS) & 5-10次 & 基本心臟切面獲取能力 \\
\hline
血管超音波 & 8-12次 & 動脈+靜脈檢查技術 \\
\hline
心導管觀摩 & 6-10台 & 涵蓋不同術式類型 \\
\hline
門診病例觀摩 & 50-80例 & 涵蓋主要心血管疾病譜 \\
\hline
\end{tabular}
\end{table}

\subsubsection{每日記錄表標準化}
建立標準化的每日記錄表格：

\begin{table}[h]
\centering
\caption{每日量化記錄表格式（新增）}
\begin{tabular}{|p{3cm}|p{2cm}|p{2cm}|p{5cm}|}
\hline
\textbf{日期} & \textbf{項目} & \textbf{數量} & \textbf{備註/學習重點} \\
\hline
\multirow{4}{*}{yyyy/mm/dd} & 病患接診 & X位 & 病名、診斷、學習重點 \\
\cline{2-4}
& POCUS & X次 & 切面種類、操作難度 \\
\cline{2-4}
& 血管超音波 & X次 & 檢查部位、異常發現 \\
\cline{2-4}
& 心導管觀摩 & X台 & 術式類型、複雜程度 \\
\hline
\multicolumn{2}{|c|}{\textbf{累計進度}} & \multicolumn{2}{|c|}{與目標比較，剩餘所需} \\
\hline
\end{tabular}
\end{table}

\subsection{2. 急症經驗保障機制}

\subsubsection{當前問題}
STEMI沒跟到是運氣問題，但急症處理能力是心臟科核心技能。

\subsubsection{解決方案}
\begin{enumerate}
\item \textbf{模擬訓練強制安排}：
   \begin{itemize}
   \item 每週至少1次STEMI案例模擬
   \item 使用標準化病人或高擬真人偶
   \item 包含從診斷到治療決策完整流程
   \end{itemize}

\item \textbf{案例討論會}：
   \begin{itemize}
   \item 每週2次急性心血管案例討論
   \item 回顧近期STEMI、急性心衰竭案例
   \item 學生主導案例報告與治療討論
   \end{itemize}

\item \textbf{急診科協作}：
   \begin{itemize}
   \item 安排半天急診科輪值
   \item 專注於胸痛病患評估
   \item 學習急症鑑別診斷流程
   \end{itemize}
\end{enumerate}

\subsection{3. 評分系統標準化}

\subsubsection{量化指標納入評分}
將量化指標完成度納入評分系統：

\begin{table}[h]
\centering
\caption{修訂後評分權重配置}
\begin{tabular}{|p{4cm}|p{2cm}|p{6cm}|}
\hline
\textbf{評分項目} & \textbf{權重} & \textbf{評分內容} \\
\hline
EPA基礎能力 & 60\% & 10項EPAs的Level評估 \\
\hline
量化指標達成 & 20\% & 基於表格化完成度評分 \\
\hline
學習態度表現 & 15\% & 記錄完整性、積極性、專業度 \\
\hline
特殊貢獻加分 & 5\% & 創新學習、額外成果 \\
\hline
\end{tabular}
\end{table}

\subsubsection{分級標準}
\begin{table}[h]
\centering
\caption{總評成績分級標準}
\begin{tabular}{|c|c|p{8cm}|}
\hline
\textbf{分數} & \textbf{等級} & \textbf{標準描述} \\
\hline
90-100 & A (優良) & 超越期待，EPA多項達Level 2-3，量化指標完成優異 \\
\hline
80-89 & B (良好) & 達到期待，EPA多項達Level 1-2，量化指標完成良好 \\
\hline
70-79 & C (合格) & 基本達標，EPA多為Level 1，量化指標基本完成 \\
\hline
60-69 & D (待改進) & 部分不足，需加強指導與延長學習 \\
\hline
<60 & F (不合格) & 嚴重不足，建議重新安排完整學習計畫 \\
\hline
\end{tabular}
\end{table}

\subsection{4. 中期評估機制}

增加第1週結束的中期評估：

\begin{itemize}
\item \textbf{量化進度檢核}：完成度是否符合預期
\item \textbf{EPA初步評估}：識別需加強的技能
\item \textbf{學習困難診斷}：及早介入輔導
\item \textbf{第2週重點調整}：個人化學習計畫
\end{itemize}

\section{未來推廣建議}

\subsection{短期實施（1個月內）}
\begin{enumerate}
\item 修訂計畫書，加入量化指標系統
\item 建立標準化記錄表格範本
\item 訓練指導醫師使用EPA評估
\item 建立急症模擬訓練教材
\end{enumerate}

\subsection{中期發展（3個月內）}
\begin{enumerate}
\item 收集5-8位學生實施數據
\item 調整量化指標合理性
\item 優化評分權重配置
\item 建立案例討論資料庫
\end{enumerate}

\subsection{長期目標（6個月內）}
\begin{enumerate}
\item 開發電子化學習管理系統
\item 建立跨科室協作網絡
\item 發展VR/AR急症訓練工具
\item 發表教學成果論文
\end{enumerate}

\section{總結}

這次兩週的心臟內科見習教學實驗取得了顯著成功。周煜珉醫師的92分優良成績證明了系統化教學設計的有效性。透過：

\begin{itemize}
\item \textbf{結構化計畫書} → 確保學習系統性
\item \textbf{量化指標管理} → 提供客觀評估依據  
\item \textbf{EPA標準化評估} → 建立能力分級標準
\item \textbf{即時記錄分享} → 促進反思與改善
\end{itemize}

我們建立了一套可複製、可改善的心臟內科見習教學模式。未來透過持續改善急症訓練保障機制和量化管理系統，相信能為更多見習學生提供高品質的心臟內科學習經驗。

\vspace{1cm}
\noindent
\begin{minipage}{0.5\textwidth}
\textbf{指導醫師：}\\
謝慕揚 醫師\\
心血管中心
\end{minipage}
\begin{minipage}{0.5\textwidth}
\textbf{日期：}\\
2025年9月14日
\end{minipage}

\end{document}